\section{End-Term University Examination (60 MCQs)}

\ExamHeader{End-Term University Examination (60 MCQs)}{Set C (Model Paper 3)}{3 Hours}{60 Marks}{-0.25 Marks}

\noindent
\textbf{Official University Examination Structure:}
\begin{enumerate}[topsep=2pt,itemsep=1pt]
    \item This paper contains \textbf{60 Multiple-Choice Questions (MCQs)} carrying \textbf{1 mark each} (Total = 60 Marks).
    \item \textbf{Curriculum Weightage:} ~20 MCQs from Units 1--3 and ~40 MCQs heavily concentrated in \textbf{Unit 4} (Quantum Mechanics), \textbf{Unit 5} (Semiconductor Physics), and \textbf{Unit 6} (Engineering Materials).
    \item Each incorrect response incurs a deduction of \textbf{0.25 marks}.
\end{enumerate}
\vspace{3mm}

\subsection*{Part I: Foundational Units 1, 2, 3 Review (Questions 1 to 20)}
\mcqitem{1}{Poisson's equation for electric potential in a medium of charge density $\rho$ is:}{$\nabla^2 V = -\rho/\epsilon_0$}{$\nabla^2 V = 0$}{$\nabla V = -\rho$}{$\nabla^2 V = \mu_0 J$}
\mcqitem{2}{If a vector field satisfies $\nabla \times \vec{F} = 0$, it is:}{Irrotational}{Solenoidal}{Rotational}{Non-conservative}
\mcqitem{3}{Maxwell's displacement current density is defined as:}{$\frac{\partial \vec{D}}{\partial t}$}{$\sigma \vec{E}$}{$\vec{E} \times \vec{H}$}{$\nabla \times \vec{H}$}
\mcqitem{4}{Maxwell's second equation $\nabla \cdot \vec{B} = 0$ implies:}{No magnetic monopoles exist}{Charge conservation}{EM wave speed is c}{Ohm's law}
\mcqitem{5}{Poynting vector is given by $\vec{S} = $}{$\vec{E} \times \vec{H}$}{$\vec{E} \cdot \vec{B}$}{$\vec{E} / \vec{B}$}{$\frac{1}{2}\epsilon E^2$}
\mcqitem{6}{Stimulated emission was theoretically discovered by:}{Albert Einstein (1916)}{Theodore Maiman}{Max Planck}{Niels Bohr}
\mcqitem{7}{The wavelength of light emitted by a He-Ne laser is:}{$632.8\text{ nm}$}{$1060\text{ nm}$}{$900\text{ nm}$}{$532\text{ nm}$}
\mcqitem{8}{Nd:YAG laser operates in which spectral region?}{Near Infrared ($1060\text{ nm}$)}{Visible red}{Ultraviolet}{X-ray}
\mcqitem{9}{Hologram records which properties of light?}{Both amplitude and phase}{Amplitude only}{Phase only}{Intensity only}
\mcqitem{10}{Total internal reflection in optical fiber requires:}{$n_{\text{core}} > n_{\text{cladding}}$}{$n_{\text{core}} < n_{\text{cladding}}$}{$n_{\text{core}} = n_{\text{cladding}}$}{$n_{\text{cladding}} = 1$}
\mcqitem{11}{Numerical Aperture (NA) is calculated as:}{$\sqrt{n_1^2 - n_2^2}$}{$n_1 - n_2$}{$n_1 + n_2$}{$\sqrt{n_1/n_2}$}
\mcqitem{12}{The cutoff normalized frequency $V$ for single-mode fiber is:}{$V \le 2.405$}{$V > 2.405$}{$V = 10$}{$V = 0$}
\mcqitem{13}{The optical window with the absolute lowest attenuation in silica fiber is:}{$1550\text{ nm}$}{$850\text{ nm}$}{$1310\text{ nm}$}{$632.8\text{ nm}$}
\mcqitem{14}{In medical endoscopy, image transmission is achieved via:}{Coherent fiber bundle}{Incoherent bundle}{Single core fiber}{Prism}
\mcqitem{15}{The ratio of He to Ne in a He-Ne gas laser is:}{10:1}{1:1}{1:10}{100:1}
\mcqitem{16}{The intrinsic impedance of free space is approximately:}{$377\,\Omega$}{$50\,\Omega$}{$100\,\Omega$}{$0\,\Omega$}
\mcqitem{17}{Stokes' theorem relates surface integral of curl with:}{Line integral around boundary curve}{Volume integral}{Point value}{Scalar gradient}
\mcqitem{18}{Ruby laser is pumped by:}{Xenon flash tube}{Electric discharge}{Chemical reaction}{Thermal heat}
\mcqitem{19}{In graded-index optical fiber, the refractive index profile is:}{Parabolic (maximum at axis)}{Constant}{Linear decrease}{Step profile}
\mcqitem{20}{Einstein's relation states that $B_{12} = $}{$B_{21}$}{$A_{21}$}{$2 B_{21}$}{$0$}

\subsection*{Part II: Core Quantum, Semiconductor, and Materials Focus (Questions 21 to 60)}
\mcqitem{21}{Planck's quantum hypothesis states that energy of radiation of frequency $\nu$ is:}{$E = h\nu$}{$E = mc^2$}{$E = \frac{1}{2}h\nu$}{$E = h/\nu$}
\mcqitem{22}{In photoelectric effect, the maximum kinetic energy of ejected electrons depends on:}{Frequency of incident light}{Intensity only}{Exposure time}{Angle of incidence}
\mcqitem{23}{The de Broglie wavelength of a particle with momentum $p$ is:}{$\lambda = h/p$}{$\lambda = p/h$}{$\lambda = h p$}{$\lambda = \sqrt{h/p}$}
\mcqitem{24}{Electron diffraction in the Davisson-Germer experiment proved:}{Wave nature of electrons}{Particle nature of light}{Photoelectric effect}{Radioactivity}
\mcqitem{25}{For matter waves, the group velocity $v_g$ equals:}{Particle velocity $v$}{Phase velocity $v_p$}{Speed of light $c$}{Zero}
\mcqitem{26}{Heisenberg's uncertainty relation for position and momentum is:}{$\Delta x \Delta p \ge \hbar/2$}{$\Delta x \Delta p = 0$}{$\Delta x \Delta p \le \hbar$}{$\Delta x / \Delta p \ge \hbar$}
\mcqitem{27}{The probability density of finding a quantum particle is given by:}{$|\psi(x,t)|^2$}{$\psi(x,t)$}{$\frac{d\psi}{dx}$}{$\int \psi dx$}
\mcqitem{28}{The energy of a particle confined in a 1D box of length $L$ is proportional to:}{$n^2$}{$n$}{$1/n$}{$1/n^2$}
\mcqitem{29}{The zero-point energy ($n=1$) of a particle in a 1D box is:}{$\frac{h^2}{8mL^2}$}{0}{$\infty$}{$\frac{h^2}{2mL^2}$}
\mcqitem{30}{The Time-Independent Schrödinger Equation in 1D is:}{$-\frac{\hbar^2}{2m}\frac{d^2\psi}{dx^2} + V(x)\psi = E\psi$}{$\frac{d\psi}{dt} = E\psi$}{$\nabla^2 \psi = 0$}{$\frac{d^2\psi}{dx^2} = E\psi$}
\mcqitem{31}{The Wiedemann-Franz Law relates:}{Thermal and electrical conductivities ($K/\sigma = L T$)}{Current and voltage}{Heat and work}{Mass and energy}
\mcqitem{32}{In an intrinsic semiconductor, the Fermi level $E_F$ is located:}{In the middle of the forbidden gap}{In the conduction band}{In the valence band}{Above conduction band}
\mcqitem{33}{Direct bandgap semiconductors (GaAs) are preferred for LEDs because:}{Recombination emits photons directly without phonons}{They are cheap}{They have zero resistance}{They have indirect transitions}
\mcqitem{34}{At absolute zero ($T = 0\text{ K}$), the Fermi-Dirac probability $f(E)$ for $E < E_F$ is:}{1}{0}{0.5}{0.25}
\mcqitem{35}{At $E = E_F$ for any temperature $T > 0\text{ K}$, the probability $f(E_F)$ is:}{0.5 (50\%)}{1}{0}{0.25}
\mcqitem{36}{A negative Hall coefficient ($R_H < 0$) indicates majority carriers are:}{Electrons (N-type)}{Holes (P-type)}{Ions}{Photons}
\mcqitem{37}{The Hall coefficient $R_H$ is given by:}{$R_H = \frac{1}{n q}$}{$R_H = n q$}{$R_H = q/n$}{$R_H = 1/(n^2 q)$}
\mcqitem{38}{A solar cell generates electrical power based on the:}{Photovoltaic effect at a p-n junction}{Seebeck effect}{Peltier effect}{Hall effect}
\mcqitem{39}{In a P-type semiconductor, the majority charge carriers are:}{Holes}{Electrons}{Photons}{Phonons}
\mcqitem{40}{Doping silicon with phosphorus (pentavalent) creates a(n):}{N-type semiconductor}{P-type semiconductor}{Intrinsic semiconductor}{Insulator}
\mcqitem{41}{Superconductivity was discovered in 1911 by:}{Heike Kamerlingh Onnes}{Albert Einstein}{Max Planck}{Niels Bohr}
\mcqitem{42}{The expulsion of magnetic flux from a superconductor ($B = 0$) is called:}{Meissner Effect}{Hall Effect}{Photoelectric Effect}{Compton Effect}
\mcqitem{43}{A superconductor inside the superconducting state is a:}{Perfect diamagnet ($\chi_m = -1$)}{Paramagnet}{Ferromagnet}{Dielectric}
\mcqitem{44}{The critical magnetic field varies with temperature according to:}{$H_c(T) = H_0 [1 - (T/T_c)^2]$}{$H_c(T) = H_0 (1 - T/T_c)$}{$H_c(T) = H_0 (1 + T/T_c)$}{$H_c(T) = H_0 (T/T_c)^2$}
\mcqitem{45}{Type-I superconductors are known as:}{Soft superconductors}{Hard superconductors}{High-Tc cuprates}{Ceramics}
\mcqitem{46}{Type-II superconductors exhibit a mixed (vortex) state between:}{$H_{c1}$ and $H_{c2}$}{$0$ and $H_{c1}$}{$H_{c2}$ and $\infty$}{None}
\mcqitem{47}{Above the Curie temperature $T_c$, a ferromagnetic material becomes:}{Paramagnetic}{Diamagnetic}{Superconducting}{Dielectric}
\mcqitem{48}{Magnetic susceptibility of a diamagnetic material is:}{Small and negative}{Large and positive}{Zero}{Infinite}
\mcqitem{49}{Soft magnetic materials have:}{Narrow hysteresis loop and low coercivity}{Wide hysteresis loop}{High coercivity}{Zero permeability}
\mcqitem{50}{The direct piezoelectric effect converts:}{Mechanical stress into electrical voltage}{Electric field into strain}{Heat into light}{Current into magnetic field}
\mcqitem{51}{The converse (inverse) piezoelectric effect converts:}{Alternating electric voltage into mechanical vibration}{Stress into voltage}{Heat into current}{Light into sound}
\mcqitem{52}{SQUID magnetometer operates using:}{Superconducting Josephson junctions}{Semiconductor diodes}{Phototubes}{Thermocouples}
\mcqitem{53}{High-temperature superconductors (HTS like YBCO) operate above the boiling point of:}{Liquid Nitrogen ($77\text{ K}$)}{Liquid Helium ($4.2\text{ K}$)}{Room temperature}{Water}
\mcqitem{54}{Maglev trains utilize which superconducting property for levitation?}{Meissner effect flux expulsion}{High resistance}{Piezoelectric effect}{Hall effect}
\mcqitem{55}{The electric displacement field $\vec{D}$ is related to $\vec{E}$ and polarization $\vec{P}$ by:}{$\vec{D} = \epsilon_0 \vec{E} + \vec{P}$}{$\vec{D} = \epsilon_0 \vec{E} - \vec{P}$}{$\vec{D} = \vec{E}/\vec{P}$}{$\vec{D} = \vec{P}$}
\mcqitem{56}{Non-polar dielectrics have:}{Zero permanent dipole moment in the absence of field}{Permanent dipoles}{Infinite conductivity}{Negative permittivity}
\mcqitem{57}{Polar dielectrics (such as water) exhibit:}{Permanent electric dipole moments}{Zero dipole moment}{Superconductivity}{Ferromagnetism}
\mcqitem{58}{In the Kronig-Penney model, energy bandgaps occur at Brillouin zone boundaries where:}{$k = \pm n\pi/a$}{$k = 0$}{$k = \infty$}{$k = a/\pi$}
\mcqitem{59}{The effective mass $m^*$ of an electron in an energy band is given by:}{$m^* = \hbar^2 / \left(\frac{d^2 E}{dk^2}\right)$}{$m^* = \hbar k$}{$m^* = \frac{dE}{dk}$}{$m^* = m_0$}
\mcqitem{60}{The relaxation time $\tau$ of conduction electrons in copper is approximately:}{$10^{-14}\text{ s}$}{$1\text{ s}$}{$10^{-6}\text{ s}$}{$10^{-3}\text{ s}$}

\newpage
\subsection*{Complete Answer Key \& Technical Rationales}

\vspace{4mm}
\begin{center}
{\small\textbf{End-Term Answer Key \& Concise Explanations (Questions 1 to 60)}}
\vspace{2mm}
\footnotesize
\begin{tabularx}{\textwidth}{l c X l c X}
\toprule
\textbf{Q\#} & \textbf{Ans} & \textbf{Technical Rationale} & \textbf{Q\#} & \textbf{Ans} & \textbf{Technical Rationale} \\
\midrule
1 & \textbf{(a)} & Electrostatic potential Poisson equation. & 31 & \textbf{(a)} & $K/\sigma = L T$ in metals. \\
2 & \textbf{(a)} & Curl-free fields are irrotational. & 32 & \textbf{(a)} & $n = p = n_i \implies E_F$ in middle of gap. \\
3 & \textbf{(a)} & Time rate of change of electric displacement field. & 33 & \textbf{(a)} & Radiative recombination occurs directly across bandgap. \\
4 & \textbf{(a)} & Magnetic field lines form continuous closed loops. & 34 & \textbf{(a)} & All states below Fermi level are completely occupied. \\
5 & \textbf{(a)} & Poynting energy flux vector ($W/m^2$). & 35 & \textbf{(a)} & $f(E_F) = \frac{1}{1 + e^0} = 0.5$. \\
6 & \textbf{(a)} & Einstein derived $A$ and $B$ coefficients. & 36 & \textbf{(a)} & Negative Hall coefficient signifies N-type semiconductor. \\
7 & \textbf{(a)} & Standard red He-Ne laser emission line. & 37 & \textbf{(a)} & Definition of Hall coefficient. \\
8 & \textbf{(a)} & 1060 nm is in near infrared. & 38 & \textbf{(a)} & Light creates electron-hole pairs separated by p-n junction field. \\
9 & \textbf{(a)} & Interference between object and reference beam records 3D information. & 39 & \textbf{(a)} & Trivalent doping creates excess majority holes. \\
10 & \textbf{(a)} & Light must travel from optically denser core to rarer cladding. & 40 & \textbf{(a)} & Pentavalent donor atoms produce N-type semiconductor. \\
11 & \textbf{(a)} & Standard formula for fiber light-gathering power. & 41 & \textbf{(a)} & Discovered in liquid helium cooled mercury at 4.12 K. \\
12 & \textbf{(a)} & Cutoff for single-mode propagation. & 42 & \textbf{(a)} & Perfect diamagnetism below $T_c$. \\
13 & \textbf{(a)} & Lowest attenuation window (~0.2 dB/km). & 43 & \textbf{(a)} & $B = \mu_0(H+M) = 0 \implies \chi_m = -1$. \\
14 & \textbf{(a)} & Coherent fiber bundles preserve 2D pixel alignment. & 44 & \textbf{(a)} & Parabolic critical magnetic field equation. \\
15 & \textbf{(a)} & Optimized 10:1 gas ratio. & 45 & \textbf{(a)} & Type-I superconductors have complete Meissner effect with low $H_c$. \\
16 & \textbf{(a)} & $\eta_0 = \sqrt{\mu_0/\epsilon_0} \approx 377\,\Omega$. & 46 & \textbf{(a)} & Abrikosov vortex state exists between $H_{c1}$ and $H_{c2}$. \\
17 & \textbf{(a)} & Stokes' integral theorem. & 47 & \textbf{(a)} & Curie-Weiss law transition to paramagnetic state. \\
18 & \textbf{(a)} & Optical pumping with flash lamp. & 48 & \textbf{(a)} & Diamagnets are weakly repelled by magnetic fields. \\
19 & \textbf{(a)} & Parabolic graded profile reduces modal dispersion. & 49 & \textbf{(a)} & Ideal for transformer cores and AC machines. \\
20 & \textbf{(a)} & Equal probability for stimulated absorption and emission. & 50 & \textbf{(a)} & Used in gas lighters, strain gauges, and piezoelectric sensors. \\
21 & \textbf{(a)} & Fundamental quantum formula. & 51 & \textbf{(a)} & Used in ultrasonic transducers and sonar emitters. \\
22 & \textbf{(a)} & Einstein: $E_k = h\nu - \phi_0$ depends on frequency. & 52 & \textbf{(a)} & Measures magnetic fields as small as $10^{-14}\text{ Tesla}$. \\
23 & \textbf{(a)} & Matter wave wavelength. & 53 & \textbf{(a)} & YBCO transitions at $93\text{ K} > 77\text{ K}$. \\
24 & \textbf{(a)} & Confirmed de Broglie hypothesis experimentally. & 54 & \textbf{(a)} & Magnetic repulsion via perfect diamagnetism. \\
25 & \textbf{(a)} & $v_g = d\omega/dk = v$. & 55 & \textbf{(a)} & Constitutive dielectric relation. \\
26 & \textbf{(a)} & Quantum uncertainty principle. & 56 & \textbf{(a)} & Dipoles are induced only when external field is applied. \\
27 & \textbf{(a)} & Born statistical interpretation. & 57 & \textbf{(a)} & Asymmetric molecular charge distribution gives permanent dipoles. \\
28 & \textbf{(a)} & $E_n = \frac{n^2 h^2}{8mL^2} \propto n^2$. & 58 & \textbf{(a)} & Bragg reflection of electron waves in crystal lattice. \\
29 & \textbf{(a)} & Minimum kinetic energy in quantum confinement. & 59 & \textbf{(a)} & Determined by band curvature in $E$-$k$ diagram. \\
30 & \textbf{(a)} & Stationary state wave equation. & 60 & \textbf{(a)} & Mean time between collisions at room temperature. \\
\bottomrule
\end{tabularx}
\end{center}
\vspace{2mm}
